\documentclass[11pt,letterpaper]{article}

% ============================================================
%  INFORME PREVIO - Proyecto AIE Summer Camp 2026
%  Tutor Adaptativo de Fundamentos de Programacion
%  Autor: Cesar Oswaldo Yunga Nieves
% ============================================================

\usepackage[utf8]{inputenc}
\usepackage[T1]{fontenc}
\usepackage[spanish]{babel}
\usepackage[margin=2.5cm]{geometry}
\usepackage{graphicx}
\usepackage{xcolor}
\usepackage{listings}
\usepackage{hyperref}
\usepackage{enumitem}
\usepackage{booktabs}
\usepackage{fancyhdr}
\usepackage{titlesec}
\usepackage{float}

% ---- Colores ----
\definecolor{rojoepn}{RGB}{155,28,40}
\definecolor{grispromt}{RGB}{240,240,240}
\definecolor{codecomment}{RGB}{100,100,100}
\definecolor{codekeyword}{RGB}{30,90,160}
\definecolor{codestring}{RGB}{40,120,70}

% ---- Enlaces ----
\hypersetup{
  colorlinks=true,
  linkcolor=rojoepn,
  urlcolor=rojoepn,
  citecolor=rojoepn
}

% ---- Encabezado / pie ----
\pagestyle{fancy}
\fancyhf{}
\lhead{\small Proyecto AIE --- Summer Camp 2026}
\rhead{\small Tutor Adaptativo}
\cfoot{\thepage}
\renewcommand{\headrulewidth}{0.4pt}

% ---- Estilo de titulos ----
\titleformat{\section}{\Large\bfseries\color{rojoepn}}{\thesection.}{0.5em}{}
\titleformat{\subsection}{\large\bfseries}{\thesubsection.}{0.5em}{}

% ---- Estilo de codigo ----
\lstdefinestyle{codigo}{
  backgroundcolor=\color{grispromt},
  basicstyle=\ttfamily\small,
  breaklines=true,
  captionpos=b,
  commentstyle=\color{codecomment}\itshape,
  keywordstyle=\color{codekeyword}\bfseries,
  stringstyle=\color{codestring},
  frame=single,
  rulecolor=\color{gray},
  numbers=left,
  numberstyle=\tiny\color{gray},
  showstringspaces=false,
  tabsize=2,
  xleftmargin=1.5em
}
\lstset{style=codigo}

% ---- Caja para instrucciones de captura (borrar al insertar imagen real) ----
\newcommand{\capturaplaceholder}[1]{%
  \begin{figure}[H]
  \centering
  \fbox{\parbox{0.85\textwidth}{\centering \vspace{1.5cm}
  \textbf{\color{rojoepn}[ INSERTAR CAPTURA AQUI ]}\\[0.3cm]
  \textit{#1}\vspace{1.5cm}}}
  \end{figure}
}

% ---- Comando para insertar una captura real con pie de foto ----
\newcommand{\captura}[2]{%
  \begin{figure}[H]
  \centering
  \fbox{\includegraphics[width=0.95\textwidth]{#1}}
  \caption{#2}
  \end{figure}
}

\begin{document}

% ============================================================
%  PORTADA
% ============================================================
\begin{titlepage}
\centering
\vspace*{2cm}

{\Huge\bfseries\color{rojoepn} Tutor Adaptativo de\\ Fundamentos de Programación\par}
\vspace{0.5cm}
{\Large Informe del Proyecto\par}
\vspace{1.5cm}

{\large\textbf{Fundamentos de Inteligencia Artificial}\par}
{\large Summer Camp 2026 --- CyberMinds EPN\par}
\vspace{2cm}

\begin{tabular}{rl}
\textbf{Autor:} & César Oswaldo Yunga Nieves \\[0.3cm]
\textbf{Modalidad:} & Individual \\[0.3cm]
\textbf{Repositorio:} & \url{https://github.com/Kaiman-p/tutor-adaptativo-ia} \\[0.3cm]
\textbf{Fecha:} & Agosto 2026 \\
\end{tabular}

\vfill
{\small Documento de informe previo --- Versión 1.0\par}
\end{titlepage}

\tableofcontents
\newpage

% ============================================================
%  1. DEFINICION DEL PROBLEMA
% ============================================================
\section{Definición del Problema}

\subsection{Contexto}
Cualquiera que haya empezado a aprender a programar conoce esta
sensación: escribes tu código, le das a ``ejecutar'', y lo único que
recibes de vuelta es un ``correcto'' o un ``incorrecto''. Nada más. Si
falló, no sabes por qué; y si funcionó, tampoco sabes si de verdad
entendiste lo que hiciste o simplemente tuviste suerte copiando la
estructura de un ejemplo anterior.

Ese vacío es un problema real. Los libros de referencia de cada
lenguaje insisten, cada uno a su manera, en que aprender a programar
no se trata solo de que el programa ``corra'', sino de comprender los
mecanismos que hay debajo. Kernighan y Ritchie, en su clásico sobre C,
son famosos por enfocarse en que el estudiante entienda de verdad cómo
funcionan los punteros y la memoria, no solo en que el código
compile~\cite{kr}. Esa misma idea --- comprender, no memorizar --- es
la que motiva este proyecto.

El problema, entonces, es que las herramientas típicas de práctica te
dicen \textit{si} acertaste, pero no \textit{qué tan bien} entendiste
el concepto que estabas practicando. Y sin esa información, es fácil
caer en una práctica poco eficiente.

\subsection{Impacto}
¿Por qué importa esto? Porque los fundamentos son, precisamente, los
cimientos. Un estudiante que arrastra un hueco en, digamos, cómo
funcionan las funciones o la recursión, va a tropezar una y otra vez
más adelante, sin saber muy bien por qué. Sweigart, en su libro de
Python para principiantes, dedica capítulos enteros a construir estos
fundamentos paso a paso justamente porque son la base de todo lo que
viene después~\cite{sweigart}.

Cuando la única señal que recibes es ``pasó / no pasó'', terminas
repitiendo lo que ya dominas (porque te da confianza) y evitando, sin
darte cuenta, lo que no entiendes bien. El resultado es una falsa
sensación de progreso.

\subsection{La solución que propongo}
La idea es un tutor con forma de videojuego (un RPG por niveles) que no
se limite a decir ``correcto''. En vez de eso, que \textbf{analice tu
código} y te diga qué tan bien estás demostrando entender cada
concepto. En concreto, el sistema:

\begin{enumerate}[leftmargin=*]
  \item Organiza el aprendizaje en ``mundos'' (un lenguaje cada uno) y
  ``niveles'' (un concepto cada uno), y a cada nivel le acompaña una
  lección pensada para alguien que parte de cero.
  \item Usa un modelo de inteligencia artificial (un LLM) como cerebro
  del sistema: lee tu código y evalúa si \textit{demuestra}
  comprensión, no solo si da el resultado correcto.
  \item Lleva la cuenta, con un modelo probabilístico, de qué tan bien
  dominas cada concepto, y va ajustando esa estimación con cada
  intento.
  \item Decide solo qué hacer contigo: avanzarte, hacerte repetir,
  generarte un ejercicio de refuerzo, o explicarte el concepto de otra
  forma si ves que no lo captaste.
\end{enumerate}

% ============================================================
%  2. RESULTADO DE IA
% ============================================================
\newpage
\section{Qué hace la IA en este proyecto}

Según el alcance del Summer Camp, la IA tenía que ser el motor real de
la solución, no un adorno. En este proyecto la IA cumple, de hecho, con
\textbf{dos} de los resultados válidos al mismo tiempo:

\begin{itemize}[leftmargin=*]
  \item \textbf{Personalización:} el contenido y el ritmo se adaptan a
  cada persona. Si dominas algo, avanzas; si te cuesta, el sistema se
  detiene, te refuerza, o te explica el concepto de otra manera.
  \item \textbf{Automatización:} después de cada intento, el sistema
  decide solo qué sigue --- sin que nadie tenga que revisarlo a mano ---
  y hasta genera ejercicios nuevos cuando hace falta.
\end{itemize}

Vale la pena aclarar por qué esto no es ``un chatbot pegado a un
formulario''. La diferencia está en qué se hace con la respuesta de la
IA: en vez de mostrarte un texto para que lo leas, el sistema le pide a
la IA que analice tu código y devuelva un resultado \textit{con
estructura} (un JSON con puntajes), y ese resultado alimenta
directamente a otro modelo que toma las decisiones. Es decir, la salida
de la IA no es el producto final que ves --- es un dato interno que
mueve toda la maquinaria. Por eso decimos que la IA es el núcleo
computacional, y no un extra.

% ============================================================
%  3. ARQUITECTURA
% ============================================================
\section{Arquitectura de Implementación}

\subsection{Cómo fluyen los datos}
Cada vez que envías una solución, pasa lo siguiente, en orden:

\begin{enumerate}[leftmargin=*]
  \item Escribes tu código para el ejercicio y le das a enviar.
  \item El \textbf{sandbox} lo ejecuta de verdad --- compila C con
  \texttt{gcc}, Java con \texttt{javac}, y corre las consultas SQL
  sobre una base de datos SQLite real --- para ver si funciona.
  \item El \textbf{LLM} (puede ser Groq, Anthropic o Gemini, se pueden
  intercambiar) lee tu código y clasifica qué tan bien dominas cada
  concepto, devolviéndolo como JSON.
  \item El \textbf{modelo BKT} toma ese resultado y actualiza su
  estimación de cuánto dominas el concepto.
  \item El \textbf{recomendador} mira esa estimación y decide qué
  sigue: avanzarte, hacerte repetir, generarte un ejercicio de refuerzo
  o reforzar el concepto más adelante.
\end{enumerate}

\captura{img/arquitectura.png}{Arquitectura del motor de IA: flujo desde el envío del estudiante hasta la decisión del sistema.}

\subsection{Los cuatro lenguajes y sus referencias}
Elegí cuatro lenguajes a propósito, porque cada uno representa una forma
distinta de pensar la programación, y entre todos obligan al sistema a
demostrar que funciona en escenarios muy diferentes:

\begin{itemize}[leftmargin=*]
  \item \textbf{Python} --- dinámico y flexible, ideal para empezar. La
  currícula de este mundo se apoya en el enfoque para principiantes de
  Sweigart~\cite{sweigart}, que construye los fundamentos sin dar nada
  por sentado.
  \item \textbf{C} --- de bajo nivel, con manejo manual de memoria y
  punteros. Aquí la referencia es el clásico de Kernighan y
  Ritchie~\cite{kr}, escrito por los propios creadores del lenguaje.
  \item \textbf{Java} --- orientado a objetos y de tipado estático. Para
  este mundo tomé como guía el enfoque didáctico de \textit{Head First
  Java}~\cite{headfirst}, pensado para que los conceptos de POO
  ``entren'' de verdad.
  \item \textbf{SQL} --- declarativo, distinto a todos los demás: no
  describes \textit{cómo} hacer algo, sino \textit{qué} quieres
  obtener. La referencia es el libro de Beaulieu~\cite{beaulieu}, un
  estándar para aprender consultas desde cero.
\end{itemize}

\subsection{Componentes técnicos}
\begin{table}[H]
\centering
\small
\begin{tabular}{@{}p{4cm}p{9cm}@{}}
\toprule
\textbf{Módulo} & \textbf{Responsabilidad} \\
\midrule
\texttt{content/*.json} & Currícula: teoría, lecciones, ejercicios,
tests y soluciones (32 niveles). \\
\texttt{sandbox.py} & Ejecuta y verifica la solución del estudiante,
con una estrategia por lenguaje. \\
\texttt{llm\_judge.py} & \textbf{Núcleo de IA.} Llama al LLM para
clasificar el dominio conceptual en JSON. \\
\texttt{mastery.py} & Modelo BKT: mantiene P(dominio) por concepto. \\
\texttt{recommender.py} & Motor de decisión sobre la siguiente
acción. \\
\texttt{exercise\_generator.py} & Genera ejercicios remediales con IA. \\
\texttt{app.py} & Interfaz visual (Streamlit). \\
\bottomrule
\end{tabular}
\caption{Módulos principales del sistema.}
\end{table}

\subsection{Herramientas y tecnologías utilizadas}
\begin{itemize}[leftmargin=*]
  \item \textbf{Lenguaje base:} Python 3.
  \item \textbf{Interfaz:} Streamlit.
  \item \textbf{Modelo de dominio:} Bayesian Knowledge Tracing (BKT),
  implementación propia.
  \item \textbf{Proveedores de LLM (intercambiables):} Groq, Anthropic
  y Google AI Studio (Gemini).
  \item \textbf{Ejecución de código:} \texttt{gcc} (C),
  \texttt{javac}/\texttt{java} (Java), \texttt{sqlite3} (SQL),
  intérprete de Python.
  \item \textbf{Control de versiones:} Git y GitHub.
\end{itemize}

% ============================================================
%  4. SISTEMA DE APRENDIZAJE (3 AYUDAS)
% ============================================================
\section{Sistema de Aprendizaje Progresivo}

Más allá de evaluar, el sistema está diseñado para \textbf{enseñar} a
alguien sin conocimiento previo. Cada nivel ofrece tres niveles de
ayuda, de menor a mayor grado de asistencia, para que el estudiante
recurra a cada uno solo cuando lo necesite. Dos de estas tres ayudas
usan la IA en tiempo real, reforzando el criterio de
\textbf{Personalización}.

\subsection{Ayuda 1: Lección completa (contenido estático)}
Cada uno de los 32 niveles incluye una lección escrita para
principiantes totales: una frase-resumen inicial (``En una frase:
\ldots''), una explicación con analogía de la vida real, y un ejemplo
resuelto \textit{distinto} al ejercicio a resolver (para no revelar la
respuesta).

\captura{img/leccion.png}{Lección completa desplegada: frase-resumen (``En una frase\ldots''), explicación con analogía y ejemplo resuelto, todo para principiantes sin conocimiento previo.}

\subsection{Ayuda 2: Pista paso a paso (IA en tiempo real)}
Si el estudiante está atascado, un botón solicita a la IA una pista de
razonamiento --- que guía sin revelar el código. El prompt de sistema
que fuerza este comportamiento es:

\begin{lstlisting}[language=Python, caption={Prompt de la pista (src/llm\_judge.py)}]
HINT_SYSTEM_PROMPT = """Eres un tutor de programacion ayudando a un
estudiante principiante que esta atascado en un ejercicio. Te daran el
enunciado del ejercicio y el codigo base (sin resolver).

Tu tarea es dar una PISTA en pasos de razonamiento -- NUNCA escribas el
codigo completo de la solucion ni la linea exacta que hay que escribir.
Divide el problema en 3 a 5 pasos de pensamiento (que revisar primero,
que estructura de control conviene usar, que caso especial no hay que
olvidar), en lenguaje simple para alguien sin experiencia previa.

Responde en espanol, en texto plano, con los pasos numerados, sin
codigo completo."""
\end{lstlisting}

\captura{img/pista.png}{Ejecución real del prompt de la pista: la IA (Gemini) genera pasos de razonamiento para el ejercicio, sin escribir el código de la solución.}

\subsection{Ayuda 3: Solución paso a paso (solucionario detallado)}
Como último recurso, el estudiante puede abrir la solución completa del
ejercicio, explicada \textbf{símbolo por símbolo} (qué hace cada
\texttt{\&}, cada \texttt{;}, cada comilla, cada \texttt{return}). Está
pensada para depurar el propio intento cuando falla por un detalle
mínimo de sintaxis, no para copiar.

\captura{img/solucion.png}{Solución paso a paso desplegada: el código de la solución seguido de la explicación detallada, símbolo por símbolo.}

\subsection{Re-explicación adaptativa (IA en tiempo real)}
Adicionalmente, un botón permite pedirle a la IA que re-explique el
concepto de una forma \textit{completamente distinta} si la explicación
original no se entendió. El prompt correspondiente:

\begin{lstlisting}[language=Python, caption={Prompt de re-explicación (src/llm\_judge.py)}]
REEXPLAIN_SYSTEM_PROMPT = """Eres un profesor paciente que explica
fundamentos de programacion a alguien que nunca ha programado antes.
Se te dara un concepto y su explicacion actual, que la persona NO
entendio. Tu trabajo es explicarlo de nuevo, de una forma COMPLETAMENTE
DISTINTA a la explicacion original -- usa una analogia diferente, un
ejemplo de la vida cotidiana distinto, y un lenguaje todavia mas
simple. Maximo 3 parrafos cortos. Responde en espanol, en texto plano."""
\end{lstlisting}

\captura{img/reexplicacion.png}{Ejecución real del prompt de re-explicación: la IA explica el mismo concepto con una analogía completamente distinta (notas adhesivas en un refrigerador).}

Este sistema de tres ayudas fue una decisión de diseño posterior a la
versión inicial: tras probar el prototipo, se identificó que un tutor
que solo evalúa no sirve para \textit{aprender} desde cero, así que se
añadió contenido pedagógico progresivo y dos usos adicionales de IA en
tiempo real (pista y re-explicación).

% ============================================================
%  5. EVIDENCIA: EJECUCION Y PROMPTS
% ============================================================
\section{Evidencia de Ejecución}

\subsection{El programa en ejecución}
La siguiente captura muestra la aplicación arrancando, con el mapa de
niveles de un mundo (lenguaje) y la barra lateral de configuración.

\captura{img/arranque.png}{La aplicación en ejecución: mapa de niveles del mundo seleccionado y barra lateral de configuración (proveedor de IA, modelo, usuario).}

\captura{img/inicio_c.png}{Vista inicial de un mundo (C): la teoría del nivel, la lección desplegable, los botones de ayuda y el enunciado del ejercicio, tal como los ve el estudiante al empezar.}

\subsection{Verificación de los 4 lenguajes con análisis real de IA}
Cada uno de los cuatro lenguajes fue probado enviando una solución y
confirmando dos cosas: que los tests pasan (``Tests superados: 100\%'')
y que la retroalimentación proviene de la IA en tiempo real (etiqueta
``fuente: llm''). Esto verifica de extremo a extremo tanto el
\textbf{sandbox} (que compila/ejecuta el código real) como el
\textbf{núcleo de IA}.

\captura{img/py_llm.png}{PYTHON: ejercicio resuelto, tests superados y retroalimentación con ``fuente: llm''.}

\captura{img/c_llm.png}{C: ejercicio resuelto y verificado (compilado con gcc real), con retroalimentación de la IA.}

\captura{img/java_llm.png}{JAVA: ejercicio resuelto y verificado (compilado con javac real), con retroalimentación de la IA.}

\captura{img/sql_llm.png}{SQL: consulta resuelta y verificada (ejecutada sobre SQLite real), con retroalimentación de la IA.}

\subsection{El prompt de sistema (núcleo de IA)}
El siguiente es el mensaje fijo que se envía al LLM antes del código
del estudiante. Obliga al modelo a responder únicamente en formato
JSON estructurado, lo que permite integrar su salida directamente en
la lógica del programa (a diferencia de un chatbot que respondería en
texto libre).

\begin{lstlisting}[language=Python, caption={Prompt de sistema del evaluador (src/llm\_judge.py)}]
JUDGE_SYSTEM_PROMPT = """Eres un evaluador tecnico de codigo para
una plataforma educativa de fundamentos de programacion. Recibiras:
el enunciado de un ejercicio, los conceptos que evalua, y el codigo
que un estudiante escribio como solucion.

Tu tarea es analizar el CODIGO EN SI (no solo si el resultado es
correcto) y devolver UNICAMENTE un JSON valido con esta forma exacta:

{
  "concept_scores": {"<concepto>": <float 0.0-1.0>},
  "detected_issues": ["<descripcion de un error detectado>"],
  "feedback": "<1-2 frases de retroalimentacion, en espanol>"
}
"""
\end{lstlisting}

% ============================================================
%  5. USO DE IA COMO COPILOTO
% ============================================================
\section{Uso de la IA como Copiloto (Trazabilidad)}

Conforme a los requisitos de documentación del Summer Camp, se
registra el uso de IA como copiloto durante el desarrollo.

\subsection{Distinción importante}
Deben distinguirse dos usos de IA en este proyecto:
\begin{enumerate}[leftmargin=*]
  \item \textbf{IA como copiloto de desarrollo:} la herramienta que
  ayudó a \textit{construir} el proyecto (esta sección). Es análoga a
  usar un editor de código o un compilador.
  \item \textbf{IA en tiempo de ejecución:} la que la \textit{propia
  aplicación} usa cada vez que un estudiante resuelve un ejercicio
  (Groq/Gemini/Anthropic). Esta es el núcleo computacional exigido por
  la rúbrica.
\end{enumerate}

\subsection{Herramienta de copiloto utilizada}
\begin{table}[H]
\centering
\small
\begin{tabular}{@{}p{4cm}p{9cm}@{}}
\toprule
\textbf{Herramienta} & \textbf{Uso} \\
\midrule
Claude (Anthropic) & Diseño de arquitectura, generación del esqueleto
de código (motor BKT, sandbox, integración LLM, recomendador),
currícula de los 32 niveles, depuración de la conexión con proveedores
de IA, y apoyo en la configuración de Git/GitHub. \\
OpenAI Tokenizer & Herramienta de análisis de tokens vista en clase;
usada como apoyo conceptual para entender el consumo de tokens y el
comportamiento de los modelos ``razonadores''. \\
\bottomrule
\end{tabular}
\caption{Herramientas de IA usadas como copiloto de desarrollo.}
\end{table}

\subsection{Modelos usados en tiempo de ejecución}
El sistema soporta tres proveedores intercambiables. Durante las
pruebas se verificó su funcionamiento con Google AI Studio (Gemini,
modelo \texttt{gemini-3.6-flash}), confirmando análisis conceptual
real y específico sobre el código del estudiante.

\subsection{Nota sobre depuración de proveedores}
Durante el desarrollo se probaron los tres proveedores con resultados
distintos, documentados como parte del proceso de ingeniería:
\begin{itemize}[leftmargin=*]
  \item \textbf{Groq:} la conexión externa funcionaba (verificado con
  \texttt{curl}), pero fallaba de forma intermitente dentro de la
  aplicación. Se identificó que el modelo \texttt{gpt-oss-120b} es un
  modelo ``razonador'' que consume parte del límite de tokens en
  razonamiento interno; se ajustó \texttt{max\_tokens} en
  consecuencia.
  \item \textbf{Gemini:} funcionando y verificado \textit{end-to-end}.
\end{itemize}
El diseño con proveedor intercambiable permitió que un fallo de
conectividad con un proveedor no bloqueara el proyecto.

% ============================================================
%  COMO EJECUTAR
% ============================================================
\section{Cómo Ejecutar el Proyecto}

El proyecto es multiplataforma (Windows, macOS o Linux). No requiere
una máquina virtual --- esta se usó únicamente en el entorno de
desarrollo del autor.

\subsection{Requisitos}
\begin{itemize}[leftmargin=*]
  \item Python 3 instalado.
  \item Una API key de un proveedor de IA (Groq, Anthropic o Google AI
  Studio; todos tienen capa gratuita).
  \item Para el mundo C: el compilador \texttt{gcc}. Para el mundo
  Java: el JDK completo (\texttt{javac}). Python y SQL no requieren
  nada adicional.
\end{itemize}

\subsection{Instalación y ejecución}
\begin{lstlisting}[language=bash, caption={Pasos para ejecutar el proyecto}]
# 1. Clonar el repositorio
git clone https://github.com/Kaiman-p/tutor-adaptativo-ia.git
cd tutor-adaptativo-ia

# 2. Instalar dependencias
pip install -r requirements.txt

# 3. Ejecutar la aplicacion
streamlit run src/app.py
\end{lstlisting}

La aplicación se abre en el navegador (\texttt{localhost:8501}). En la
barra lateral se configura el proveedor de IA, el modelo y la API key.
Sin API key, el sistema funciona en modo \textit{fallback} (solo
verificación de tests, sin análisis conceptual de IA).

% ============================================================
%  6. RESULTADOS Y VERIFICACION
% ============================================================
\section{Resultados y Verificación}

Los 32 niveles (4 lenguajes $\times$ 8 niveles) fueron verificados con
ejecución real, no solo revisados manualmente:
\begin{itemize}[leftmargin=*]
  \item \textbf{Python, C, SQL:} las soluciones de referencia se
  probaron contra los tests reales de cada nivel mediante un script
  automatizado.
  \item \textbf{Java:} compilado y ejecutado con \texttt{javac} y
  \texttt{java} reales en la máquina del autor.
  \item \textbf{Interfaz completa:} probada con el framework oficial
  de testing de Streamlit (\texttt{AppTest}), simulando interacciones
  reales.
\end{itemize}

% ============================================================
%  LIMITACIONES Y TRABAJO FUTURO
% ============================================================
\section{Limitaciones y Trabajo Futuro}

Como todo prototipo desarrollado en un tiempo acotado, el sistema tiene
limitaciones conocidas que se documentan con honestidad:

\begin{itemize}[leftmargin=*]
  \item \textbf{No detecta copia directa de la solución.} El sistema
  mide corrección y comprensión conceptual (vía el LLM), pero no puede
  distinguir si el estudiante razonó la respuesta o copió la solución
  mostrada en la ayuda ``Ver solución paso a paso''. El diseño confía
  en un uso progresivo de las ayudas, como cualquier libro de texto con
  soluciones al final.
  \item \textbf{Es monousuario.} El progreso se guarda por nombre de
  usuario en archivos locales; no hay un sistema de cuentas real. Dos
  personas con el mismo nombre compartirían progreso.
  \item \textbf{Dependencia de un servicio externo de IA.} El análisis
  conceptual requiere conexión a un proveedor de LLM. Sin conexión, el
  sistema se degrada al modo \textit{fallback} (funcional, pero sin
  análisis conceptual).
  \item \textbf{Cobertura de lenguajes.} Se implementaron 4 lenguajes;
  extender a más requiere crear la currícula correspondiente (el motor
  ya está diseñado para soportarlo sin cambios estructurales).
\end{itemize}

\subsection{Trabajo futuro}
\begin{itemize}[leftmargin=*]
  \item Extender a más lenguajes sobre el mismo motor genérico.
  \item Añadir un sistema de cuentas real para uso multiusuario.
  \item Incorporar una interfaz visual más rica (animaciones, sonidos)
  al estilo de un videojuego completo.
  \item Explorar detección de patrones de copia para reforzar la
  integridad del aprendizaje.
\end{itemize}

% ============================================================
%  7. CONCLUSIONES
% ============================================================
\section{Conclusiones}

El proyecto demuestra que la IA puede ser el núcleo computacional real
de una herramienta educativa, no un componente decorativo. La
combinación de un LLM (para análisis conceptual del código) con un
modelo probabilístico clásico (BKT, para modelar el dominio) permite
tomar decisiones pedagógicas autónomas y personalizadas, cumpliendo
simultáneamente los resultados de Personalización y Automatización
definidos en el alcance del Summer Camp.

El diseño modular y con proveedor de IA intercambiable demostró su
valor durante el desarrollo, al permitir sortear problemas de
conectividad sin rediseñar el sistema.

% ============================================================
%  DECLARACION DE USO DE IA
% ============================================================
\vspace{1cm}
\noindent\rule{\textwidth}{0.4pt}

\subsection*{Declaración de uso de Inteligencia Artificial}
Este proyecto fue desarrollado con el apoyo de \textbf{Claude
(Anthropic)} como herramienta de copiloto de programación, conforme a
los lineamientos del Summer Camp que permiten y exigen documentar el
uso de IA. El autor revisó, probó y comprende cada componente del
sistema, y tomó las decisiones de diseño del proyecto (selección de
lenguajes, elección del modelo de dominio, proveedores de IA, y
enfoque pedagógico). El detalle completo de las sesiones de trabajo
con IA se encuentra en el archivo \texttt{docs/registro\_prompts.md}
del repositorio.

% ============================================================
%  BIBLIOGRAFIA
% ============================================================
\newpage
\begin{thebibliography}{9}

\bibitem{kr}
Kernighan, B. W. \& Ritchie, D. M. (1988).
\textit{The C Programming Language} (2\textsuperscript{a} ed.).
Prentice Hall. ISBN 978-0131103627. Disponible en:
\url{https://seriouscomputerist.atariverse.com/media/pdf/book/C\%20Programming\%20Language\%20-\%202nd\%20Edition\%20(OCR).pdf}.

\bibitem{sweigart}
Sweigart, A. (2020).
\textit{Automate the Boring Stuff with Python} (2\textsuperscript{a} ed.).
No Starch Press. Disponible en línea:
\url{https://automatetheboringstuff.com/}.

\bibitem{headfirst}
Sierra, K. \& Bates, B. (2005).
\textit{Head First Java} (2\textsuperscript{a} ed.).
O'Reilly Media. ISBN 978-0596009205. Disponible en:
\url{https://archive.org/details/HeadFirstJava2ndEdition_201511}.

\bibitem{beaulieu}
Beaulieu, A. (2020).
\textit{Learning SQL} (3\textsuperscript{a} ed.).
O'Reilly Media. ISBN 978-1492057611. Disponible en:
\url{https://archive.org/details/learningsql0000beau_s2h5}.

\end{thebibliography}

\end{document}
